
\medskip

%Cet exercice est constitué de 5 questions indépendantes.
%
%\medskip

\begin{enumerate}
\item Sur la figure ci-dessous, chacun des quadrilatères quad1, quad2 et quad3 est l'image du quadrilatère TRAP par une transformation.

\begin{center}
\psset{unit=1cm}
\begin{pspicture}(0,-0.4)(7,8)
%\psgrid
\def\quad{\pspolygon(1.8,0)(0.8,1.4)(0,1.4)(0,0)\psframe(0.15,0.15)\psframe(0,1.4)(0.15,1.25)}%PART
\rput(0.5,4){\quad}
\uput[dl](0.5,4){\small P}\uput[ur](2.3,3.6){\small A}\uput[ur](1.3,5.4){\small R}\uput[ul](0.5,5.4){\small T}
\psline(1,0)(7,6)\uput[dr](3.6,2.6){D}\uput[dr](6,5){E}
\psdots[dotstyle=+,dotangle=45](3.6,2.6)(6,5)\rput{90}(2.3,2.2){\quad}
\rput{-80}(5.4,-0.2){\psscalebox{-1 1}{\quad}}\rput(5.9,0.6){quad1}
\rput(2.9,6.4){\quad}\rput(3.6,6.9){quad2}\rput(1.7,3){quad3}
\end{pspicture}
\end{center}

%Recopier les trois phrases ci-dessous sur la copie et compléter, sans justifier, chacune d'elles par le numéro de l'une des transformations proposées dans le tableau qui suit:
	\begin{enumerate}
		\item Le quadrilatère quad1 est l'image du quadrilatère TRAP par la transformation numéro 6.
		\item Le quadrilatère quad2 est l'image du quadrilatère TRAP par la transformation numéro 1. 
		\item Le quadrilatère quad3 est l'image du quadrilatère TRAP par la transformation numéro 2.
	\end{enumerate}

%\medskip
%	
%\begin{tabularx}{\linewidth}{|X|X|}\hline
%\textbf{Transformation numéro 1} : translation qui transforme le point D en le point E.&
%\textbf{Transformation numéro 4} : translation qui transforme le point E en le point D.\\ \hline
%\textbf{Transformation numéro 2} : rotation de centre A et d'angle 90\degres{} dans le sens contraire des aiguilles d'une montre.&
%\textbf{Transformation numéro 5} : rotation de centre A et d'angle 120\degres{} dans le sens contraire des aiguilles d'une montre.\\ \hline
%\textbf{Transformation numéro 3} : symétrie centrale de centre D.&
%\textbf{Transformation numéro 6} : symétrie axiale d'axe (DE).\\ \hline
%\end{tabularx}
\item %Développer et réduire l'expression suivante: 

%\[(2x - 3)(- 5 + 2x) - 4 + 6x.\]
$(2x - 3)(- 5 + 2x) - 4 + 6x = - 10x + 4x^2  + 15 - 6x - 4 + 6x = 4x^2 - 10x + 11$.
\item %Résoudre l'équation suivante : 

%\[(x - 6)(5x - 2) = 0.\]
Un produit de facteurs est nul si l'un des facteurs est nul, donc :

$(x - 6)(5x - 2) = 0$ si $x - 6 = 0$ ou $5x - 2 = 0$ soit :

$x = 6$ ou $5x = 2$ et enfin $x = 6$ ou $x = \dfrac{2}{5} = \dfrac{4}{10} = 0,4$.

$S = \{0,4~;~6\}$.
\item 
	\begin{enumerate}
		\item %Décomposer, sans justifier, en produit de facteurs premiers les nombres \np{1386} et \np{1716}.
\starredbullet~~$\np{1386} = 9 \times 154 = 9 \times 14 \times 11 = 2 \times 3^2 \times 7 \times 11$ ;

\starredbullet~~$\np{1716} = 6 \times 286 = 6 \times 2 \times 143 = 6 \times 2 \times 13 \times 11 = 2^2 \times 3 \times 11 \times 13$.
		\item %En déduire la forme irréductible de la fraction : $\dfrac{\np{1386}}{\np{1716}}$.
		$\dfrac{\np{1386}}{\np{1716}} = \dfrac{2 \times 3^2 \times 7 \times 11}{2^2 \times 3 \times 11 \times 13} = \dfrac{3 \times 7}{2 \times 13} = \dfrac{21}{26}$.
	\end{enumerate}
\item %Les coordonnées géographiques de la ville appelée Jokkmokk sont environ: 67\degres{} Nord et 19\degres{} Est. 

%Placer approximativement la ville de Jokkmokk sur le planisphère.

\begin{center}
\begin{pspicture}(-6,-4)(6,4)


\rput(0,0){\includegraphics[width=12cm]{mappemonde.eps}}
\psdots[dotstyle=+,dotangle=45,dotscale=2](0.3,2.3)
\uput[d](0.3,2.3){Jokkmokk}
\end{pspicture}
\end{center}

\end{enumerate}

\bigskip

